\section{Discussion}
\label{sec:discussion}

\subsection{Limitations}
\label{sec:disc-limits}

The pre-registered ablation in Section~\ref{sec:experimental-plan} isolates
a single dimension of the recursive design space: the depth-limit value
$\depthlim$ on a fixed serialization, fixed action-targeting scheme, and
fixed (frozen) LLM encoder. Other dimensions of the recursive design,
such as branching factor at each node, adaptive depth conditioned on
problem difficulty, and per-verb depth budgets, are not exercised in this
paper's plan. We also do not test alternative tree encoders against the
serialization-plus-frozen-LLM approach argued for in
Section~\ref{sec:method-serialization}. A natural follow-up would freeze
the framework paper's heads and swap only the state encoder between (a)
flat serialization with depth markers and (b) a Tree-LSTM
\citep{tai2015improved} or graph attention network
\citep{velickovic2018gat} operating on the parent-pointer graph, leaving
$\Qpsi$ and $\ptheta$ otherwise unchanged. We flag this as a natural
extension and do not pursue it here.

The productive-versus-unproductive distinction is operational rather than
syntactic: it is defined in terms of whether terminal verifier reward
reaches the leaves of the induced tree within the depth budget
(Definition~\ref{def:productive}). This operational definition has the
methodological virtue of avoiding ad hoc syntactic classifications, but
it also creates marginal cases. A hypothesis that is provably equivalent
to its parent (a tautological refinement) produces a structurally legal
recursive descent that may, depending on the verifier's behavior on the
duplicated claim, classify either way. Likewise, in domains with a
learnable verifier such as a process reward model \citep{lightman2024verify},
the productive-versus-unproductive distinction would inherit any noise
the learned verifier introduces; the framework paper avoids this by
restricting Bellman targets to non-learned verifier signals only, and our
analysis here preserves that restriction. The cleanest interpretation of
the depth-limit mechanism is in domains where the verifier is exact
(SymPy normalization, formal proof kernels, unit-test execution) and the
productive-versus-unproductive boundary is sharp.

The depth limit itself is a hard mask on the available action set (the
removal of \textsc{Decompose}, \textsc{CheckCase}, and depth-incrementing
\textsc{FormHypothesis} from $\Aset(s)$ at the limit), not a learned
attention. An alternative would be to learn the recursion structure
end-to-end, treating the decision to recur as a policy parameter rather
than as an MDP-level constraint. We did not pursue this path because it
would couple the recursion question with capacity questions and complicate
the ablation: a learned-recursion variant that underperforms could be
failing either because end-to-end recursion is wrong or because the
specific learning machinery does not yet have enough capacity. The
hard-mask design isolates the structural-prior claim cleanly, at the cost
of giving up adaptive flexibility within a single trajectory. Finally,
this is a methodological companion to the framework paper; empirical
results from the ablation grid are deferred to that paper's
implementation. The methodological contribution stands independent of
empirical confirmation, but the productive-versus-unproductive claim is
falsifiable and we make it so via the pre-registered hypotheses in
Section~\ref{sec:exp-ablation}.

\subsection{Extensions}
\label{sec:disc-extensions}

The fixed $\depthlim \in \{3, 4\}$ default and the global-mask ablation
both admit natural relaxations. The simplest is a \emph{learned per-problem
depth limit}: a small head reads the problem statement and predicts the
useful tree depth, which the MDP then enforces via the same hard-mask
mechanism. Slightly more sophisticated versions condition on running
search statistics, annealing the limit as confidence in a subtree's
hypothesis grows. Both extensions are compatible with the verifier-only
reward principle and with the hard-mask treatment, since the prediction
is part of the environment specification rather than of the reward
signal. A second extension is \emph{per-verb recursion budgets}: some
verbs (e.g., \textsc{Decompose}) may merit deeper recursion than others
(e.g., \textsc{CheckCase}), and the availability table in
Section~\ref{sec:method-depthlimit} can be amended per-verb without
changing the policy or value architecture. Both extensions can be
plugged into the ablation grid once the headline result is in.

A third extension worth surfacing is \emph{soft penalties on recursion}
as an alternative to the hard mask. We argued in
Section~\ref{sec:method-depthlimit} that soft penalties violate the
framework paper's verifier-only-reward principle and admit unbounded
simulation descent. Both critiques stand on first principles. Empirically,
however, a small soft penalty may still be a useful comparator: if the
hard mask and a comparable soft penalty produce similar empirical
behavior, the architectural cost of the hard mask is harder to justify;
if they differ, the difference is informative about which mechanism the
$Q$-learning system relies on. We flag this as worth empirical check in
follow-up work without endorsing it as the recommended default.

A fourth extension is \emph{cross-domain recursion patterns}. The
productive-versus-unproductive distinction was developed for
verifier-anchored math, where the verifier accepts intermediate calls and
the natural decomposition depth is small. How does the distinction
transfer to code generation, where the verifier is a unit-test suite that
typically fires only at the global solution? To formal proof, where the
verifier (Lean kernel) accepts only fully discharged sub-claims? To
open-ended deliberation, where there is no verifier at all? We expect
the productive regime to shrink and the depth-limit's importance to grow
as the verifier becomes more global and as the natural decomposition
depth becomes more variable. The cross-domain question is the obvious
next paper after the present one's intra-domain ablation lands.

Finally, a DAG generalization (multi-parent hypotheses) and a cross-trajectory
caching mechanism (verified sub-claims re-used as facts in later problems)
are both compatible with the implicit-hierarchy stance and both worth
exploring. The schema's parent-pointer field would need to become a list
rather than a single identifier; the depth-distance metric used by
$\depthlim$ would need a poset analogue. Neither change is fundamental,
but the canonicalization questions deferred in
Section~\ref{sec:method-canon} would need answering before either
extension is empirically tractable.

\subsection{When Recursion Pays Off}
\label{sec:disc-when}

We have argued throughout that the productive-versus-unproductive
distinction is operational, not syntactic. The same architectural
question can be asked at the level of problem classes: which classes
benefit measurably from recursive structural decomposition, and which do
not? The pattern we predict, based on the verifier-reach analysis of
Section~\ref{sec:method-productive}, has two features. First, the
problem class must admit \emph{natural verifier-checkable sub-goals}:
the kind of decomposition where a tentative parent claim can be split
into sub-claims whose conjunction implies the original, and each
sub-claim is itself a verifier-checkable object. Competition mathematics
\citep{hendrycks2021math}, multi-claim formal proof \citep{deepmind2024alphaproof,
yang2023leandojo}, and code generation with unit-test decomposition
\citep{li2022alphacode} all have this property by construction. Second,
the verifier must accept intermediate calls: a verifier that scores
sub-claims independently of the parent gives the productive regime its
reward signal at the leaves of the implicit tree, where the $Q$-function
needs Bellman targets most.

Problems missing either feature gain little from recursion. A
one-shot answer problem with a global verifier (the verifier only fires
on the terminal commit) yields no leaf-level signal beneath; the depth
budget that recursion costs is not offset by intermediate verifier hits,
and $\Qpsi$ has no informative Bellman target for the interior nodes.
Open-ended deliberation, where there is no verifier at all, falls in the
same category: recursive decomposition is structurally available, but
nothing makes the productive mode preferable to the unproductive one
from the agent's perspective. Within the math-and-proof setting that
motivates the framework paper, we expect the productive regime to exist
in a measurable range; outside that setting, the depth-limit mechanism
may still be useful as a defense against rumination but the
expressivity benefit will be smaller.

\subsection{Threats Specific to the Recursion Claim}
\label{sec:disc-threats}

The experimental plan's threats subsection (Section~\ref{sec:exp-threats})
covers the broad threats inherited from the framework paper: compute
leaks, verifier coverage, benchmark contamination, seed variance. We
note here the threats specific to the depth-limit claim. The ablation
may show no difference between R0, R3, and R$\infty$, in which case the
productive-versus-unproductive distinction is real but the depth-limit
mechanism is the wrong way to operationalize it; an alternative such as
per-verb budgets or learned recursion may then become the next experiment.
The ablation may show R3 dominating both endpoints but only by a small
margin that does not justify tracking parent-pointer state at all; in
that case the framework paper's deferral of recursion to a phase-2
forward-compatibility hook is the correct long-term answer rather than
a temporary one. The ablation's three settings ($\depthlim \in \{0, 3,
\infty\}$) may also fail to span the truly interesting regime. We
mitigate this by running a finer-grained sweep on a held-out subset
after the headline three-setting result is in
(Section~\ref{sec:exp-threats}).

Two further threats deserve specific mention. Verifier coverage at
intermediate depths can inflate or deflate the productive-recursion rate
artificially: a verifier that accepts a wider class of intermediate
claims will register more leaves as productive even if the underlying
deliberation is unchanged, and the empirical depth distribution we
report in Section~\ref{sec:exp-secondary} must be read with the
verifier's coverage profile in mind. Cross-seed variance, finally, is
likely to be high in the $\depthlim = \infty$ condition specifically:
unbounded recursion makes the trajectory length distribution
heavy-tailed, and long trajectories carry more accumulated variance into
the $Q$-learning targets. We pre-commit to reporting mean and standard
deviation across at least three training seeds and to flagging any
condition whose seed variance exceeds the across-condition effect.

\subsection{Implicit versus Explicit Hierarchy: A Methodological Note}
\label{sec:disc-implicit}

The hierarchical reinforcement learning literature has spent two decades
architecting explicit hierarchy. The options framework
\citep{sutton1999between} introduces sub-policies with their own
termination conditions; MAXQ value decomposition
\citep{dietterich2000hierarchical} introduces sub-task value functions
with cross-coupling terms; feudal networks \citep{vezhnevets2017feudal}
introduce manager modules operating at coarser timescales than worker
modules. Each of these designs makes hierarchy a property of the policy
or value architecture: separate networks, separate parameter sets,
separate optimization signals at each level. The framework paper adopts
the opposite methodological stance, and the present companion paper
develops it in depth. Hierarchy here is a property of the \emph{state},
not of the policy: $\Qpsi$ is a flat categorical head over the
seven-verb action set, $\ptheta$ is a flat policy over the same set, the
frozen LLM is a single shared encoder, and recursion enters entirely
through the parent-pointer field in the state schema and the depth
integer attached to each entity.

The trade-off worth surfacing is this. Explicit hierarchy gives the
agent a structural inductive bias that constrains learning: the manager
cannot choose worker-level primitives directly, the option must run to
its termination condition, the sub-task value function decomposes as a
sum over levels. The bias accelerates learning when the designer's
chosen decomposition matches the task's natural decomposition; it
distorts learning when it does not. Implicit hierarchy lets the agent
discover (or fail to discover) the appropriate hierarchical decomposition
from the data, at the cost of denying it the inductive bias that
explicit machinery would have supplied. For pre-trained language models
already carrying world knowledge about how mathematical claims and
sub-claims relate, we expect the implicit approach to extract more
transfer than an explicit hierarchy whose decomposition rules are baked
in by the system designer. The LLM's pretraining already supports
structured text reasoning; adding explicit hierarchical machinery on top
of it would duplicate that capacity, with the redundancy paid for in
architectural complexity and in fine-tuning signal that has to choose
between the implicit and explicit hierarchy whenever they disagree.

The commitment is, then, methodological: keep the architecture as simple
as possible (flat categorical $\Qpsi$, flat policy $\ptheta$, frozen
LLM as both encoder and stochastic transition model), and let hierarchy
emerge from the state representation rather than from the policy
architecture. The recursive paper's design is the test of whether this
commitment carries through when the state itself becomes
tree-structured. The depth-limit mechanism, the identifier-based action
targeting, and the flat serialization with depth markers are all
consequences of the same underlying choice: hierarchy lives in the
state, the architecture stays flat, and the verifier signals do all the
load-bearing work of teaching the system which trajectories matter.
